\documentclass[10pt,a4paper]{amsart}
\usepackage[margin=19mm]{geometry}
\usepackage{amsmath,amssymb,mathtools}
\usepackage[scr=rsfs,cal=euler]{mathalfa}
\usepackage{xfrac}
\usepackage[unicode,hidelinks,bookmarks=false]{hyperref}
\newcommand{\ie}{i.e.\ }
\newcommand{\cf}{cf.\ }
\newcommand{\resp}{respectively\ }
\newcommand{\Subsection}{\subsection}
\providecommand{\nobreakdash}{\nobreak\hbox{-}}
\setlength{\emergencystretch}{2em}
\multlinegap0pt
\usepackage{placeins}
\usepackage{color}
\usepackage{psfrag}
\usepackage{latexsym}
\usepackage{tabularx}
\usepackage{bm}
\usepackage{lscape}
\usepackage{color}
\usepackage{url}
\usepackage{upgreek}
\usepackage{units}
\usepackage{fancyhdr}
\definecolor{OliveGreen}{rgb}{0,0.6,0}
\usepackage{enumitem}
\usepackage{ulem}
\renewcommand{\baselinestretch}{1.3}
\setlength{\parskip}{\baselineskip}
\setlength{\parindent}{0pt}%
\newtheorem{theorem}{Theorem}[section]
\newtheorem{lemma}[theorem]{Lemma}
\newtheorem{constraint}[theorem]{Constraint}
\newtheorem{corollary}[theorem]{Corollary}
\newtheorem{proposition}[theorem]{Proposition}
\newtheorem{conjecture}[theorem]{Conjecture}
\newtheorem{result}[theorem]{Result}
\newtheorem{fact}{Fact}
\newtheorem{definition}[theorem]{Definition}
\newtheorem{remark}[theorem]{Remark}
\newtheorem{example}[theorem]{Example}
\newtheorem{assumption}[theorem]{Assumption}
\newtheorem{property}[theorem]{Property}
\numberwithin{equation}{section}
\numberwithin{theorem}{section}
\numberwithin{table}{section}
\numberwithin{figure}{section}
\DeclareMathOperator*{\argmax}{arg\,max}
\DeclareMathOperator*{\argmin}{arg\,min}
\newcommand{\N}{{\rm I\!N}}
\newcommand{\R}{{\rm I\!R}}
\newcommand{\PI}{{\rm I\!\Pi}}
\newcommand{\e}{{\rm e}}
\newcommand{\E}{{\mathbb E}}
\newcommand{\p}{{\mathbb P}}
\newcommand{\V}{{\mathbb V}}
\newcommand{\q}{{\mathbb Q}}
\begin{document}
\title[Posterior and Likelihood Sensitivity in Bayesian Distributio]{Posterior and Likelihood Sensitivity in Bayesian Distributionally Robust Optimization}
\author{Gotoh, Kim, Lim [Posterior and Likelihood Sensitivity in Bayesian DRO] Posterior and Likelihood Sensitivity in; Bayesian Distributionally Robust Optimization ạte; Jun-ya Gotoh, Andrew E.B. Lim, Michael Jong Kim}
\date{}
\maketitle
\noindent\emph{Source and licence.} Posterior and Likelihood Sensitivity in Bayesian Distributionally Robust Optimization. Version arXiv:2605.31306v1 (CC BY 4.0). This material is distributed under the Creative Commons Attribution 4.0 International licence, http://creativecommons.org/licenses/by/4.0/, as recorded in the arXiv OAI licence metadata for this exact version and shown on the abstract page https://arxiv.org/abs/2605.31306v1. Full rights notice: licenses/arxiv-2605.31306-CC-BY-4.0.txt.\par\smallskip
\noindent\textit{Editorial scope.} Counted unit: the original \texttt{assumption} environment \texttt{ass:Bayes} at source lines 266--290 together with its complete proof at lines 320--340; the delivered document keeps source lines 197--561 so that every referenced label and every citation resolves inside it. Native editable source is the paper's own concatenated TeX; the AMS wrapper is editorial formatting only. No publisher class, upstream script or unrelated artwork is redistributed.
\par\smallskip\noindent\textit{Reference note.} This article ships no compiled bibliography (biblatex); the entries delivered here are composed from the author's own BibTeX (.bib) fields, with no field invented and no entry dropped.
\section{Worst-case sensitivity} 
\label{sec:WCS}

We use ideas from \cite{gotoh2026sensitivity} to define worst-case sensitivity with respect to the posterior and the likelihood.  


\subsection{Uncertainty in the posterior and likelihood}

We begin by defining deviations from the posterior $\rho({\rm d}\theta)$ and likelihood ${\mathbb P}({\rm d}y | \theta)$. 

Let $\eta$ be a distribution  on $\Theta$  and ${\rm d}( \eta | \rho)$ be a deviation measure of $\eta$ from $\rho$. For every $\theta\in\Theta$ let ${\mathbb Q}^\theta$ be an alternative to $L^\theta$ for the distribution of  $Y$ and ${\rm d}({\mathbb Q}^\theta| L^\theta)$ a measure of deviation of ${\mathbb Q}^\theta$ from $L^\theta$. It will occasionally be convenient to write ${\mathcal L} = \{ L^\theta| \theta\in\Theta\}$ when we talk about the entire family of distributions.

%Let ${\mathcal Q} = \{{\mathbb Q}^\theta\vert\theta\in\Theta\}$ be a family of probability distributions for $Y$ parameterized by $\theta\in\Theta$ and ${\mathcal L} = \{ L^\theta| \theta\in\Theta\}$ be the nominal likelihood function. For every $\theta$, ${\rm d}({\mathbb Q}^\theta| L^\theta)$ is the deviation of ${\mathbb Q}^\theta$ from $L^\theta$. 
%We define the deviation of ${\mathcal Q}$ from $\mathcal L$ as the expected deviation of the probability measure ${\mathbb Q}^\theta$ from ${\mathbb P}^\theta$ under the posterior $\rho$:
%\begin{eqnarray}
%{\rm d}({\mathcal Q}| {\mathcal L}) & = & {\mathbb E}_\mathbb \rho \big\{{\rm d}({\mathbb Q}^\theta| L^\theta)\big\}.
%\label{eq:dQL}
%\end{eqnarray}



%\subsection{Constraint model}
Let
\begin{eqnarray}
 \label{eq:uncertainty-set}
{\mathcal P}(\varepsilon) & = & \big\{ \eta\big\vert {\rm d}(\eta|\rho)\leq \varepsilon\big\},\; \\  [5pt]
{\mathcal L}^\theta(\delta) & =& \big\{{\mathbb Q}\big\vert {\rm d}({\mathbb Q}|L^\theta)\leq\delta\big\} \nonumber
\end{eqnarray}
be uncertainty sets for the posterior and likelihood distributions. We use the same notation ${\rm d}(\cdot|\cdot)$ for the deviation measure in each uncertainty set. However, different deviation measures can be used for the posterior and likelihood (see Section \ref{sec:mix}).

Consider the worst-case expected cost
\begin{eqnarray}
V(\epsilon, \delta) 
%& = & \min_{\eta \in {\mathcal P}_\phi(\varepsilon)} \min_{\{{\mathbb Q}^\theta\} \in {\mathcal L}_\phi(\delta)} {\mathbb E}_\eta {\mathbb E}_{Q^\theta} [f(x, Y)] \\
= \max_{\eta \in {\mathcal P}(\varepsilon)} {\mathbb E}_\eta \Big\{\max_{{\mathbb Q}^\theta \in {\mathcal L}^\theta(\delta)} {\mathbb E}_{Q^\theta} [f(\theta, Y)] \Big\}.
\label{eq:wc_Bayes}
\end{eqnarray}
The inner expectation is conditional on $\theta$; we consider worst-case deviations ${\mathbb Q}^\theta$ from the nominal likelihood $L^\theta$ for each $\theta$. The outer expectation considers worst-case deviations $\eta$ from the posterior $\rho$ so changes the weights on $\theta$. 

\subsection{Regularization}

Consider first the inner problem. Since the worst-case expected cost is increasing in the size of the uncertainty set $\delta$, we can write for every $\theta\in\Theta$
\begin{eqnarray}
\Psi(\theta) = \max_{{\mathbb Q}^\theta \in {\mathcal L}^\theta(\delta)} {\mathbb E}_{{\mathbb Q}^\theta} [f(\theta, Y)] 
   =    {\mathbb E}_{L^\theta}[f(\theta, Y)]  + {\mathcal A}_{L^\theta}\big(\delta; f(\theta, Y)\big)
   \label{eq:inner}
   \end{eqnarray}
where   ${\mathcal A}_{L^\theta}\big(\delta; f(\theta, Y)\big)$, the ambiguity cost, is non-negative and non-decreasing in $\delta$. 

Similarly, we can write the outer problem
\begin{eqnarray}
\max_{\eta \in {\mathcal P}(\varepsilon)} {\mathbb E}_\eta \big\{\Psi(\theta) \big\} & = & {\mathbb E}_\rho \big\{\Psi(\theta) \big\} + {\mathcal A}_{\rho}\big(\varepsilon; \Psi(\theta)\big)
\label{eq:outer}
\end{eqnarray}
where the ambiguity cost ${\mathcal A}_{\rho}\big(\varepsilon; \Psi(\theta)\big)$ is again non-negative and increasing in $\varepsilon$.


Substituting \eqref{eq:inner} into \eqref{eq:outer} it  follows that the worst-case expected cost \ref{eq:wc_Bayes} can be written
\begin{eqnarray*}
V(\epsilon, \delta)  = {\mathbb E}_\rho \Big\{{\mathbb E}_{L^\theta}[f(\theta, Y)]  \Big\} + {\mathcal A}_{\rho,{\mathcal L}}\big(\varepsilon, \delta; f(\theta, Y)\big)
\end{eqnarray*}
with ambiguity cost
\begin{eqnarray*}
{\mathcal A}_{\rho,{\mathcal L}}\big(\varepsilon, \delta; f(\theta, Y)\big) := {\mathbb E}_\rho \Big\{ {\mathcal A}_{L^\theta}\big(\delta; f(\theta, Y)\big) \Big\} + {\mathcal A}_{\rho}\Big(\varepsilon; {\mathbb E}_{L^\theta}[f(\theta, Y)]  + {\mathcal A}_{L^\theta}\big(\delta; f(\theta, Y)\big)\Big).
\end{eqnarray*}



We make the following assumptions about the ambiguity costs. The functions $g(\varepsilon)$ and $g(\delta)$ for the two ambiguity costs need not be the same.

\begin{assumption} \label{ass:Bayes}
The ambiguity costs ${\mathcal A}_{L^\theta}\big(\delta; f(\theta, Y)\big)$  and ${\mathcal A}_{\rho}\big(\varepsilon; \Psi(\theta)\big)$ are such that 
\begin{enumerate}
\item there is a non-decreasing function $g(\delta)$ such that $g(\delta)\rightarrow 0$ when $\delta\rightarrow 0$ and 
\begin{eqnarray*}
%\Psi^\delta(\theta, f(\theta, \cdot)) & =& 
%\max_{{\mathbb Q}^\theta \in {\mathcal L}(\delta)} {\mathbb E}_{Q^\theta} [f(\theta, Y)] 
 %  & = &   {\mathbb E}_{L^\theta}[f(\theta, Y)]  + 
{\mathcal A}_{L^\theta}\big(\delta; f(\theta, Y)\big) %\nonumber \\ [5pt]
   =  g(\delta) {\mathcal S}_{L^\theta}(f(\theta, Y))+ o(g(\delta))
%\label{eq:Psi}
\end{eqnarray*}
for every $\theta\in{\Theta}$;
%where 
%\begin{eqnarray*}
\item there is a non-decreasing function $g(\varepsilon)$ such that $g(\varepsilon)\rightarrow 0$ when $\varepsilon\rightarrow 0$ and 
\begin{eqnarray*}
%\max_{\eta \in {\mathcal P}(\varepsilon)} {\mathbb E}_\eta \big\{\Psi(\theta) \big\} & = & {\mathbb E}_\rho \big\{\Psi(\theta) \big\} + 
{\mathcal A}_{\rho}\big(\varepsilon; \Psi(\theta)\big) %\\ [5pt]
=  %{\mathbb E}_\rho \big\{\Psi(\theta) \big\} + 
g(\varepsilon) {\mathcal S}_\rho(\Psi(\theta)) + o(g(\varepsilon)).
\end{eqnarray*}
\item  ${\mathcal S}_{\rho}(\Psi)$ is continuous in $\Psi$. 
\end{enumerate}
\end{assumption}

Intuitively, $\Psi(\theta)$ is the worst-case expected cost when  the uncertainty set  is $\{{\mathcal L}^\theta(\delta)\}$ and
\begin{eqnarray*}
{\mathcal A}_{L^\theta}\big(\delta; f(\theta, Y)\big)= \max_{{\mathbb Q}^\theta \in {\mathcal L}^\theta(\delta)} {\mathbb E}_{{\mathbb Q}^\theta} [f(\theta, Y)] 
   -   {\mathbb E}_{L^\theta}[f(\theta, Y)]
   \end{eqnarray*} 
is the  increase in expected cost relative to the nominal. When Assumption \ref{ass:Bayes} holds,  ${\mathcal S}_{L^\theta}(f(\theta, Y))$  can be interpreted as the  sensitivity of the expected cost under worst-case perturbations; it depends on the deviation measure ${\rm d}({\mathbb Q} | L^\theta)$ that defines the uncertainty set. For example, ${\mathcal S}_{L^\theta}(f(\theta, Y))$ is the standard deviation of the cost $\sigma_{L^\theta} (f(\theta, Y))$ when ${\rm d}({\mathbb Q}|L^\theta)$ is smooth $\phi$-divergence. Sensitivity measures induced by other measures of deviation can be found in \cite{gotoh2026sensitivity} and references cited therein (see also Section \ref{sec:mix}).
   
Likewise, ${\mathbb E}_\rho \big\{\Psi(\theta) \big\}$ is the expected value of $\Psi(\theta)$ when $\theta$ has distribution $\rho$ and
${\mathcal S}_{\rho}(\Psi)$ is the sensitivity of the expected cost ${\mathbb E}_\rho \big\{\Psi(\theta) \big\}$ under worst-case perturbations of $\rho$. 
   
Uncertainty in posterior and likelihood contribute to increases in the total expected cost
   \begin{eqnarray*}
   {\mathcal A}_{\rho,{\mathcal L}}\big(\varepsilon, \delta; f(\theta, Y)\big) = \max_{\eta \in {\mathcal P}(\varepsilon)} {\mathbb E}_\eta \Big\{\max_{{\mathbb Q}^\theta \in {\mathcal L}^\theta(\delta)} {\mathbb E}_{Q^\theta} [f(\theta, Y)] \Big\} - {\mathbb E}_\rho \Big\{{\mathbb E}_{L^\theta}[f(\theta, Y)]  \Big\}.
   \end{eqnarray*}
Our goal is to understand how the sensitivity of the total cost depends on the posterior and likelihood.

The following result gives an expansion of the value function when $\varepsilon$ and $\delta$ are small. 
%The proof can be found in Appendix \ref{App:Bayes}.
\begin{theorem} \label{prop:Bayesian expansion}
Suppose that Assumption \ref{ass:Bayes} holds. Then
\begin{eqnarray*}
V(\epsilon, \delta) & = & 
{\mathbb E}_\rho\Big\{{\mathbb E}_{L^\theta}[f(\theta, Y)] \Big\} + g(\delta){\mathbb E}_\rho\Big\{{\mathcal S}_{L^\theta}(f(\theta, Y))\Big\}  +  g(\varepsilon) {\mathcal S}_\rho\Big({\mathbb E}_{L^\theta}(f(\theta, Y))\Big) \\ [5pt] 
& & + o(g(\delta)) + o(g({\varepsilon})).
\end{eqnarray*}
\end{theorem}


%\subsection{Proof of Proposition \ref{prop:Bayesian expansion}}

\begin{proof}
It follows from \eqref{eq:inner} and Assumption \ref{ass:Bayes} that
\begin{eqnarray*}
\Psi(\theta)  &= &   {\mathbb E}_{L^\theta}[f(\theta, Y)]  + {\mathcal A}_{L^\theta}\big(\delta; f(\theta, Y)\big) \\ & = & {\mathbb E}_{L^\theta}[f(\theta, Y)]  + g(\delta) {\mathcal S}_{L^\theta}(f(\theta, Y))+ o(g(\delta)).
\end{eqnarray*}
From \eqref{eq:outer} and Assumption \ref{ass:Bayes} we have
\begin{align*}
V(\varepsilon, \delta) & =  \max_{\eta \in {\mathcal P}(\varepsilon)} {\mathbb E}_\eta \big\{\Psi(\theta) \big\} \\[8pt] 
& = \max_{\eta \in {\mathcal P}(\varepsilon)} {\mathbb E}_\eta\Big\{{\mathbb E}_{L^\theta}[f(\theta, Y)]  + g(\delta) {\mathcal S}_{L^\theta}(f(\theta, Y))+ o(g(\delta))\Big\} \\[8pt]
& = {\mathbb E}_\rho\Big\{{\mathbb E}_{L^\theta}[f(\theta, Y)]+ g(\delta) {\mathcal S}_{L^\theta}(f(\theta, Y))+ o(g(\delta)) \Big\} \\[8pt]
& \qquad + g(\varepsilon){\mathcal S}_\rho\Big({\mathbb E}_{L^\theta}[f(\theta, Y)]  + g(\delta) {\mathcal S}_{L^\theta}(f(\theta, Y))+ o(g(\delta))\Big) + o(g(\varepsilon)) \\[8pt]
& = {\mathbb E}_\rho\Big\{{\mathbb E}_{L^\theta}[f(\theta, Y)] \Big\} + g(\delta){\mathbb E}_\rho\Big\{{\mathcal S}_{L^\theta}(f(\theta, Y))\Big\}  +  g(\varepsilon) {\mathcal S}_\rho\Big({\mathbb E}_{L^\theta}(f(\theta, Y))\Big) \\[8pt]
& \qquad + g(\varepsilon) \left\{{\mathcal S}_\rho\Big({\mathbb E}_{L^\theta}[f(\theta, Y)]  + g(\delta) {\mathcal S}_{L^\theta}(f(\theta, Y)) + o(g(\delta)) \Big) - {\mathcal S}_\rho\Big({\mathbb E}_{L^\theta}(f(\theta, Y))\Big)\right\}  
\\[8pt]
& 
\qquad + o(g(\delta)) + o(g({\varepsilon})) \\[8pt]
& = {\mathbb E}_\rho\Big\{{\mathbb E}_{L^\theta}[f(\theta, Y)] \Big\} + g(\delta){\mathbb E}_\rho\Big\{{\mathcal S}_{L^\theta}(f(\theta, Y))\Big\}  +  g(\varepsilon) {\mathcal S}_\rho\Big({\mathbb E}_{L^\theta}(f(\theta, Y))\Big) \\[8pt] 
& \qquad + o(g(\delta)) + o(g({\varepsilon}))
\end{align*}
where the last equality follows from the continuity of ${\mathcal S}_\rho(\Psi)$.
\end{proof}

%\hfill$\Box$



\subsection{Worst-case sensitivity}

Given likelihood function $L^\theta(y)$ ($\theta\in\Theta$) and posterior $\rho$, 
worst-case posterior sensitivity is the sensitivity of the expected cost under worst-case deviations from the posterior defined by the uncertainty set ${\mathcal P}(\varepsilon)$:
\begin{eqnarray*}
{\mathcal S}_\rho(f)  =  \lim_{\epsilon\rightarrow 0}\frac{V(\epsilon, 0)- {\mathbb E}_\rho\Big\{{\mathbb E}_{L^\theta}[f(\theta, Y)]\Big\}}{g({\epsilon})}% \nonumber \\ [5pt]
=
{\mathcal S}_\rho\Big\{{\mathbb E}_{L^\theta}[f(\theta, Y)]\Big\}.
%\label{eq:prior-sensitivity-const}
\end{eqnarray*}

Given the nominal posterior $\rho$, likelihood  $L^\theta(y)$, and uncertainty set ${\mathcal L}^\theta(\delta)$ for deviations from the likelihood, worst-case likelihood sensitivity is the sensitivity of the expected cost with respect to worst-case deviations from the nominal likelihood
\begin{eqnarray*}
{\mathcal S}_{\mathcal L}(f)  =  \lim_{\delta\rightarrow 0}\frac{V(0, \delta)- {\mathbb E}_\rho\Big\{{\mathbb E}_{L^\theta}[f(\theta, Y)]\Big\}}{g(\delta)}% \nonumber \\ [5pt]
 =  {\mathbb E}_\rho \Big\{{\mathcal S}_{L^\theta}(f(\theta, Y))\Big\}.
%\label{eq:likelihood-sensitivity-const}
\end{eqnarray*}
Note that posterior (likelihood) sensitivity ${\mathcal S}_\rho$  depends on the deviation measure ${\rm d}(\eta|\rho)$  that defines the uncertainty set ${\mathcal P}(\varepsilon)$, and likelihood sensitivity  ${\mathcal S}_{L^\theta}$ on the deviation measure  ${\rm d}({\mathbb Q}|L^\theta)$ that controls deviations from the nominal likelihood. 





In general, ${\mathcal S}_\rho$ is a generalized measure of deviation \cite{gotoh2026sensitivity,rockafellar2006generalized} so posterior sensitivity is zero when $\rho$ is a point mass (i.e., there is no uncertainty about $\theta$) or when $g(\theta) = {\mathbb E}_{L^\theta}[f(\theta, Y)]$ is constant on the support of $\rho$. %Likelihood sensitivity is non-zero unless $f(\theta, y)$ is constant as a function of $y$ on the support of $L^\theta$. 
Likelihood sensitivity is zero only when, for $\rho$-almost every $\theta$, the sensitivity of the cost under $L^\theta$ is zero. For the deviation measures considered here, this occurs when $f(\theta, y)$ is constant in $y$ on the support of $L^\theta$.



The following example considers posterior and likelihood sensitivity when deviation measures are smooth $\phi$-divergence.

\subsection{Smooth $\phi$-divergence}
 \label{eg:phi_Bayes}
 

Suppose that $\eta$ is absolutely continuous with respect to $\rho$ and  ${\mathbb Q}^\theta$ is absolutely continuous with respect to $L^\theta$ for every $\theta\in\Theta$. Let $\frac{{\rm d} \eta}{{\rm d}\rho}$ denote the Radon--Nikodym derivative of $\eta$ with respect to $\rho$, and for every $\theta\in\Theta$, $\frac{{\rm d}Q^\theta}{{\rm d}L^\theta}$ be the Radon--Nikodym derivative of $Q^{\theta}$ with respect to $L^\theta$.


For the likelihood function,  $\phi$-divergence of ${\mathbb Q}^\theta$ with respect to $L^\theta$, for every $\theta\in{\Theta}$, is defined by
\begin{eqnarray*}
 {\rm d}({\mathbb Q}^\theta|L^\theta)= {\mathbb E}_{L^\theta} \left[\phi \left(\frac{{\rm d}{\mathbb Q}^\theta}{{\rm d} L^\theta}\right) \right].
\end{eqnarray*}
In the case of smooth $\phi$-divergence, Assumption \ref{ass:Bayes} holds and the inner problem \cite{gotoh2026sensitivity} 
\begin{eqnarray*}
\max_{{\mathbb Q}^\theta \in {\mathcal L}^\theta_{\phi}(\delta)} {\mathbb E}_{Q^\theta} [f(\theta, Y)] 
=  {\mathbb E}_{L^\theta}[f(\theta, Y)]  + \sqrt{\frac{2\delta}{\phi''(1)}} {\sigma}_{L^\theta}(f(\theta, Y))+ o(\sqrt{\delta}).
\end{eqnarray*}
For the outer problem we measure the deviation of $\eta$ from $\rho$ using $\phi$-divergence:
\begin{eqnarray*}
{\rm d}( \eta | \rho) & = & {\mathbb E}_\rho \left[\phi\left(\frac{{\rm d}\eta}{{\rm d}\rho}\right)\right]
\end{eqnarray*}
and it follows that
\begin{eqnarray*}
\max_{\eta \in {\mathcal P}_\phi(\varepsilon)} {\mathbb E}_\eta[\Psi(\theta)] 
= {\mathbb E}_\rho[\Psi(\theta)] + \sqrt{\frac{2\varepsilon}{\phi''(1)}}\sigma_\rho(\Psi(\theta)) + o(\sqrt{\varepsilon}).
\end{eqnarray*}
In particular, $g(\delta) = \sqrt{\delta}$ and 
\begin{eqnarray*} 
{\mathcal S}_{L^\theta}(f) =  \sqrt{\frac{2}{\phi''(1)}}  {\sigma}_{L^\theta}(f(\theta, Y))
\end{eqnarray*}  
for the inner problem, where ${\sigma}_{L^\theta}(f(\theta, Y)) = \sqrt{ {\mathbb V}_{L^\theta}(f(\theta, Y))}$ is the standard deviation of the cost under the distribution $L^\theta$. For the outer problem,
$g(\varepsilon) = \sqrt{\varepsilon}$ and
\begin{eqnarray*}
{\mathcal S}_\rho(\Psi(\theta))  =\sqrt{\frac{2}{\phi''(1)}}\sigma_\rho(\Psi(\theta)).
\end{eqnarray*}
Note that the regularizer in the expansion of the worst-case expected costs are given by the standard deviation because we are using smooth $\phi$-divergence as the deviation measure. Other uncertainty sets give different regularizers \cite{gotoh2026sensitivity}.


It now  follows from Theorem \ref{prop:Bayesian expansion} that the worst-case expected cost 
\begin{eqnarray*}
V(\epsilon, \delta) %& = & \min_{\eta \in {\mathcal P}_\phi(\varepsilon)} {\mathbb E}_\eta \Big\{\min_{{\mathbb Q}^\theta \in {\mathcal L}_\phi(\delta)} {\mathbb E}_{Q^\theta} [f(x, Y)] \Big\}  \label{eq:WC-Bayes-penalty} \\
& = &  {\mathbb E}_\rho\Big\{{\mathbb E}_{L^\theta}[f(\theta, Y)]\Big\}+ \sqrt{\frac{2\delta}{\phi''(1)}}{\mathbb E}_\rho \Big({\sigma}_{L^\theta}(f(\theta, Y))\Big) + \sqrt{\frac{2\varepsilon}{\phi''(1)}}
{\sigma}_\rho\Big({\mathbb E}_{L^\theta}[f(\theta, Y)]\Big) \nonumber \\
& & + o(\sqrt{\epsilon}) + o(\sqrt{\delta}).
%\label{eq:WC-Bayes-penalty}
\end{eqnarray*}
Worst-case posterior  sensitivity is
\begin{eqnarray}
{\mathcal S}_\rho(f) %& = & \lim_{\epsilon\rightarrow 0}\frac{V(\epsilon, 0)- {\mathbb E}_\rho\Big\{{\mathbb E}_{L^\theta}[f(x, Y)]\Big\}}{\sqrt{\epsilon}} \nonumber \\ 
%&  = & 
%\sqrt{\frac{2}{\phi''(1)}}
= \sqrt{\frac{2}{\phi''(1)}}{\sigma}_\rho\Big({\mathbb E}_{L^\theta}[f(\theta, Y)]\Big)
\label{eq:prior-sensitivity-const}
\end{eqnarray}
and worst-case likelihood sensitivity is %worst-case sensitivity with respect to the likelihood is defined by
\begin{eqnarray}
{\mathcal S}_{\mathcal L}(f) %& = & \lim_{\epsilon\rightarrow 0}\frac{V(0, \delta)- {\mathbb E}_\rho\Big\{{\mathbb E}_{L^\theta}[f(x, Y)]\Big\}}{\sqrt{\delta}} \nonumber \\ 
%&  = &
= \sqrt{\frac{2}{\phi''(1)}} {\mathbb E}_\rho \Big\{{\sigma}_{L^\theta}(f(\theta, Y))\Big\}.
\label{eq:likelihood-sensitivity-const}
\end{eqnarray}
Both sensitivity measures depend on the posterior and objective function, but quite differently. Posterior sensitivity  is the standard deviation of the conditional expectation $g(\theta) = {\mathbb E}_{L^\theta}[f(\theta, Y)]$ and will be small if $g(\theta)$ is ``flat" on the support of $\rho$, even if the standard deviation of $\theta$ is large. 
It vanishes when the posterior variance shrinks (e.g.) because of learning. Likelihood sensitivity depends on the variance of the cost given $\theta$. It typically will not vanish even if $\theta$ is known because there is still uncertainty about the distribution $L^\theta$ of $Y$.


The paper \cite{shapiro2023bayesian} derives worst-case likelihood sensitivity  with a Kullback--Leibler uncertainty set (the inner ``max" in \eqref{eq:wc_Bayes}) but not the posterior. We complement this by introducing posterior sensitivity and by showing how the worst-case objective can be approximated by  expected cost, posterior sensitivity, and likelihood sensitivity.


\subsection{Mix and match}
\label{sec:mix}

As noted in the discussion following \eqref{eq:uncertainty-set}, there is no reason to restrict ourselves to smooth $\phi$-divergence for the deviation measure or to have the same deviation measure for the posterior and likelihood. For example, if instead of smooth $\phi$-divergence, we use 
\begin{eqnarray*}
\phi(z)=\delta_{[1-\varepsilon,\frac{1-(1-\varepsilon)\alpha}{1-\alpha}]}(z)
\end{eqnarray*} 
to define the uncertainty set for the likelihood function, we have \cite{gotoh2026sensitivity}\footnote{Given $\theta$, $\mathrm{CVaR}_{L^\theta,\alpha}(f(\theta, Y))$ is the conditional-value-at-risk of the cost $f(\theta, Y)$ at level $\alpha$ when $Y$ has distribution $L^\theta$: \begin{eqnarray*}\mathrm{CVaR}_{L^\theta,\alpha}(f(\theta, Y)):=\min_\gamma\Big\{\gamma+\frac{1}{1-\alpha}\mathbb{E}_{L^\theta}\Big(\max\big\{f(\theta, Y)-\gamma,0\big\}\Big)\Big\}\end{eqnarray*}}
\begin{eqnarray*} 
{\mathcal S}_{L^\theta}(f) = \mathrm{CVaR}_{L^\theta,\alpha}(f(\theta, Y))-\mathbb{E}_{L^\theta} (f(\theta, Y)).
\end{eqnarray*} 
(WCS for other deviation measures are also derived in \cite{gotoh2026sensitivity}). We choose this uncertainty set and deviation measure if we are concerned about misspecification of the right tail of the cost distribution. 
If smooth $\phi$-divergence is used for the posterior and this alternative is used for the likelihood, we have
\begin{align*}
{\mathcal S}_{\rho}(f) &= \sqrt{\frac{2}{\phi''(1)}}{\sigma}_\rho\Big({\mathbb E}_{L^\theta}[f(\theta, Y)]\Big), \\[5pt]
{\mathcal S}_{\mathcal L}(f) & = {\mathbb E}_\rho \Big\{\mathrm{CVaR}_{L^\theta, \alpha}(f(\theta, Y))-\mathbb{E}_{L^\theta} (f(\theta, Y))\Big\}.
\end{align*}



%Posterior sensitivity is important because it quantifies the sensitivity of the expected cost to worst-case deviations from the posterior which depends on the user-specified prior distribution. 

%\subsection{Other uncertainty sets}

%An important class of deviation measures are of the form
%\begin{eqnarray*}
%{\rm d}({\mathbb Q}|{\mathbb P}) =  {\mathbb E}_{\mathbb P}\Big[\phi\Big(\frac{d {\mathbb Q}}{d {\mathbb P}}\Big)\Big].
%\end{eqnarray*}
%Another important deviation measure is the Wasserstein distance. By selecting ${\mathbb P} \equiv \rho$ or ${\mathbb P} \equiv L^\theta$ these deviation measures can be used to form uncertainty sets posterior and likelihood sensitivity. 
%Worst-case sensitivity for different choices of $\phi$ are provided in Table \ref{table:summary} as well as for the Wasserstein metric. Although smooth $\phi$-divergence  is used for both the posterior and likelihood uncertainty sets in Section \ref{eg:phi_Bayes}, other deviation measures can be used and it is not necessary for the deviation measure for the posterior uncertainty set to be the same as that of the likelihood (as noted in the discussion immediately following \eqref{eq:uncertainty-set}).


%\begin{table}
%\begin{center}
%\caption{Worst-case sensitivities considered in Section \ref{sec:WCS}}
%\label{table:summary}
%\begin{tabular}{ll|c}
%\multicolumn{3}{c}{$\langle$i$\rangle$~ $\max_{\mathsf{q}}\{\mathbb{E}_{\mathsf{q}}(\mathsf{f})\vert d(\mathsf{q}\vert\mathsf{p})\leq \varepsilon\}$}\\
%\hline
%\multicolumn{2}{l|}{Type of divergence/worst-case objective} & Worst-case sensitivity \eqref{eq:sensitivity-general2} \\
%\hline
%\multicolumn{2}{l|}{(a) smooth $\phi$-divergence
%\qquad
%$\phi''(1)>0=\phi(1)=\phi'(1)$} & $\sqrt{\frac{2{\mathbb V}_{\mathsf{p}}(\mathsf{f})}{\phi''(1)}}$ \\ [2pt]
%(b) total variation &
%$\phi(z)=|z-1|$ & $\frac{1}{2}\big(\max(\mathsf{f})-\min(\mathsf{f})\big)$\\ [2pt]
%(c) ``budgeted'' 
%& $\phi=\delta_{[0,1+\varepsilon]}$ & $\mathbb{E}_{\mathsf{p}}(\mathsf{f})-\min(\mathsf{f})$ \\[2pt]
%(d) $(1-\varepsilon)$``mean% cost
%''$+\varepsilon$``$\alpha$-CVaR'' & 
%$\phi=\delta_{[1-\varepsilon,\frac{1-(1-\varepsilon)\alpha}{1-\alpha}]}$
%& $\mathrm{CVaR}_{\mathsf{p},
%\alpha}(\mathsf{f})-\mathbb{E}_{\mathsf{p}}
%(\mathsf{f})$\\ [5pt]
%(d') $(1-\varepsilon)$``mean
%''$+\varepsilon$``max
%'' 
%& $\phi=\delta_{[1-\varepsilon,
%\infty)}$
% & $\max(\mathsf{f})-\mathbb{E}_{\mathsf{p}}(\mathsf{f})$\\[2pt]
%(d'') symmetric ($\alpha=\frac{1}{2-\varepsilon}$) & 
%$\phi=\delta_{[1-\varepsilon,\frac{1}{1-\varepsilon}]}$ & $\mathrm{CVaR}_{\mathsf{p},\frac{1}{2}}(\mathsf{f})-\mathbb{E}_{\mathsf{p}}(\mathsf{f})$ \\[5pt]
%\multicolumn{2}{l|}{(e) type-1 Wasserstein distance 
%~w/ Lipschitz smooth $f$ 
%} & $\max\limits_{i=1,\cdots,n} \max\limits_{z_i}\frac{f(z_i)-f(Y_i)}{\|z_i-Y_i\|}$ \\
%\hline
%\end{tabular}\\
%\end{center}
%\end{table}



\section{Distributionally Robust Optimization}
\label{sec:DRO}

\subsection{Performance--sensitivity tradeoffs in DRO}
\label{sec:DROtradeoff}

Theorem \ref{prop:Bayesian expansion} shows that the  worst-case expected cost can be locally expanded as
\begin{eqnarray}
\lefteqn{\min_x \max_{\eta \in {\mathcal P}(\varepsilon)} {\mathbb E}_\eta \Big\{\max_{{\mathbb Q}^\theta \in {\mathcal L}^\theta(\delta)} {\mathbb E}_{Q^\theta} [f(x, \theta, Y)] \Big\}}
\label{eq:DRO-regularized}
\\ [5pt]
& = & \min_x {\mathbb E}_\rho\Big\{{\mathbb E}_{L^\theta}[f(x, \theta, Y)]\Big\}+ g(\delta) {\mathbb E}_\rho \Big({\mathcal S}_{L^\theta}(f(x, \theta, Y))\Big) + g(\varepsilon)  {\mathcal S}_\rho\Big({\mathbb E}_{L^\theta}[f(x, \theta, Y)]\Big) 
\nonumber  \\ [5pt] 
& & + o(g({\epsilon})) + o(g({\delta})).
\nonumber
\end{eqnarray}
This representation is related to results showing the relationship between DRO and a regularized nominal problem. For empirical optimization problems, \cite{gotoh2026sensitivity} shows that the regularization term is not just a mathematical object that connects the nominal and worst-case problems, but  has the physical interpretation as worst-case sensitivity. \eqref{eq:DRO-regularized} shows an analogous result for a Bayesian DRO problem, where the regularizer is now a weighted sum of posterior  sensitivity ${\mathcal S}_\rho({\mathbb E}_{L^\theta}[f(x, \theta, Y)])$   and likelihood sensitivity ${\mathbb E}_\rho ({\mathcal S}_{L^\theta}(f(x, \theta, Y)))$, and varying the uncertainty set sizes maps out a tradeoff between performance (expected cost) and robustness (posterior and likelihood sensitivity). 

%It is now interesting to understand how a robust DM makes tradeoffs between expected cost and posterior and likelihood sensitivity.





%\begin{remark}
%Worst-case prior and likelihood sensitivities for the penalty-version of the Bayesian DRO problem with a smooth $\phi$-divergence penalty is given in Appendix \ref{App:BayesPenalty}.
%\end{remark}







\section{Experiment}
\label{sec:experiments}


Demand $D(p)$ is normal with standard deviation $\sigma$ and expected value ${\mathbb E}[D(p)] = a + b p$ which depends on the price $p$. We assume that $\sigma$ is known and constant whereas $a$ and $b$ are constant but unknown to the DM. At the time of the decision, the decision maker has a posterior on $(a, b)$ which we assume to be normal. (This could have been obtained by updating an initial prior using data and Bayes' rule.) The DM's revenue when price is $p$ and observed demand is $D(p)$ is $p\cdot D(p)$. The expected revenue under the posterior is ${\mathbb E}_\rho [p \, D(p)] = p \cdot ({\mathbb E}_\rho[a]  + {\mathbb E}_\rho[b]\cdot p)$; the nominal DM chooses price to maximize this quantity. 

We assume a modified $\chi^2$--uncertainty set for the posterior and likelihood, i.e., $\phi(z) = \frac{1}{2}(z-1)^2$. It follows that posterior and likelihood sensitivities are given by
\begin{align}
\label{eq:pricing-senitivities}
{\mathcal S}_\rho[p\, D(p)] & = \sqrt{2} p \, \sigma_\rho [a + b p], \\
{\mathcal S}_{\mathcal L}[p\, D(p)] & = \sqrt{2}\sigma p. \nonumber
\end{align}
For this experiment, under the baseline posterior, $a$ has mean $5$ and standard deviation $0.75$, and $b$ has mean $-0.5$ and standard deviation $2.5$; we assume $a$ and $b$ are independent under the posterior. We set $\sigma = 3$.  For the nominal problem the optimal price is $p =\$5$; the optimal expected reward is $\$12.50$.


\begin{thebibliography}{99}
\bibitem[]{gotoh2026sensitivity} Gotoh, Jun-ya; Kim, Michael Jong; Lim, Andrew E.B.. Robustness Measures in Distributionally Robust Optimization. ar$\chi$iv:2507.11350 (2026)
\bibitem[]{rockafellar2006generalized} Rockafellar, R. Tyrrell; Uryasev, Stan; Zabarankin, Michael. Generalized deviations in risk analysis. Finance and Stochastics 10 (2006) 51--74
\bibitem[]{shapiro2023bayesian} Shapiro, Alexander; Zhou, Enlu; Lin, Yifan. Bayesian distributionally robust optimization. SIAM Journal on Optimization 33 (2023) 1279--1304
\end{thebibliography}
\end{document}
